\documentclass[11pt,oneside]{article}
\usepackage[a4paper,margin=27mm,headheight=15pt]{geometry}
\usepackage[T1]{fontenc}
\usepackage[utf8]{inputenc}
\IfFileExists{libertinus.sty}{\usepackage{libertinus}}{\usepackage{lmodern}}
\usepackage{microtype}
\usepackage{amsmath,amssymb,amsthm,mathtools,mathrsfs,bm}
\usepackage{booktabs,longtable,array,tabularx}
\usepackage{graphicx}
\usepackage{xcolor}
\usepackage{enumitem}
\usepackage{fancyhdr}
\usepackage{titlesec}
\usepackage{listings}
\usepackage{xurl}
\usepackage[hypertexnames=false]{hyperref}
\usepackage[nameinlink,noabbrev]{cleveref}
\usepackage{caption}
\usepackage{float}
\usepackage{placeins}

\definecolor{linkblue}{RGB}{25,75,135}
\definecolor{shadegray}{RGB}{246,247,249}
\definecolor{rulegray}{RGB}{195,201,208}
\definecolor{sourceblue}{RGB}{35,91,147}
\definecolor{newgreen}{RGB}{25,112,75}
\definecolor{conjecturered}{RGB}{145,55,55}
\hypersetup{
  colorlinks=true,
  linkcolor=linkblue,
  citecolor=linkblue,
  urlcolor=linkblue,
  pdftitle={Atomic Functions Beyond the Critical Dyadic Case},
  pdfsubject={Translation and expansion of Rvachev atomic-function material with new fractal and asymptotic results},
  pdfkeywords={Rvachev up-function, Fabius function, atomic functions, Thue-Morse, sinc product, fractal string, Lambert W, Bell polynomials, Bernoulli numbers}
}
\pagestyle{fancy}
\fancyhf{}
\fancyhead[L]{\small Atomic functions beyond the dyadic case}
\fancyhead[R]{\small\nouppercase{\leftmark}}
\fancyfoot[C]{\thepage}
\renewcommand{\headrulewidth}{0.4pt}
\titleformat{\section}{\Large\bfseries}{\thesection}{0.7em}{}
\titleformat{\subsection}{\large\bfseries}{\thesubsection}{0.7em}{}
\titleformat{\subsubsection}{\normalsize\bfseries}{\thesubsubsection}{0.7em}{}
\setcounter{secnumdepth}{3}
\setcounter{tocdepth}{2}
\setlist[itemize]{leftmargin=2em,itemsep=0.25em,topsep=0.4em}
\setlist[enumerate]{leftmargin=2.3em,itemsep=0.25em,topsep=0.4em}
\setlength{\emergencystretch}{2em}

\newtheorem{theorem}{Theorem}[section]
\newtheorem{proposition}[theorem]{Proposition}
\newtheorem{lemma}[theorem]{Lemma}
\newtheorem{corollary}[theorem]{Corollary}
\newtheorem{conjecture}[theorem]{Conjecture}
\theoremstyle{definition}
\newtheorem{definition}[theorem]{Definition}
\newtheorem{algorithm}[theorem]{Algorithm}
\newtheorem{example}[theorem]{Example}
\theoremstyle{remark}
\newtheorem{remark}[theorem]{Remark}
\newtheorem{warning}[theorem]{Warning}

\crefname{theorem}{theorem}{theorems}
\Crefname{theorem}{Theorem}{Theorems}
\crefname{proposition}{proposition}{propositions}
\Crefname{proposition}{Proposition}{Propositions}
\crefname{lemma}{lemma}{lemmas}
\Crefname{lemma}{Lemma}{Lemmas}
\crefname{corollary}{corollary}{corollaries}
\Crefname{corollary}{Corollary}{Corollaries}
\crefname{conjecture}{conjecture}{conjectures}
\Crefname{conjecture}{Conjecture}{Conjectures}
\crefname{definition}{definition}{definitions}
\Crefname{definition}{Definition}{Definitions}
\crefname{algorithm}{algorithm}{algorithms}
\Crefname{algorithm}{Algorithm}{Algorithms}
\crefname{example}{example}{examples}
\Crefname{example}{Example}{Examples}
\crefname{remark}{remark}{remarks}
\Crefname{remark}{Remark}{Remarks}
\crefname{warning}{warning}{warnings}
\Crefname{warning}{Warning}{Warnings}

\newcommand{\R}{\mathbb R}
\newcommand{\C}{\mathbb C}
\newcommand{\Q}{\mathbb Q}
\newcommand{\N}{\mathbb N}
\newcommand{\Z}{\mathbb Z}
\newcommand{\Up}{\operatorname{up}}
\newcommand{\sinc}{\operatorname{sinc}}
\newcommand{\e}{\mathrm e}
\newcommand{\ii}{\mathrm i}
\newcommand{\dd}{\,\mathrm d}
\newcommand{\Var}{\operatorname{Var}}
\newcommand{\E}{\mathbb E}
\newcommand{\Prob}{\mathbb P}
\newcommand{\bigO}{\mathcal O}
\newcommand{\smallo}{o}
\newcommand{\supp}{\operatorname{supp}}
\newcommand{\dist}{\operatorname{dist}}
\newcommand{\wt}{\operatorname{wt}_2}
\newcommand{\diam}{\operatorname{diam}}
\newcommand{\TV}{d_{\mathrm{TV}}}
\newcommand{\Wass}{W_{\infty}}
\newcommand{\eps}{\varepsilon}
\newcommand{\floor}[1]{\left\lfloor #1\right\rfloor}
\newcommand{\fract}[1]{\left\{#1\right\}}
\newcommand{\repo}{\nolinkurl{Analysis/FabiusFunction/docs}}
\newcommand{\Smark}{\textcolor{sourceblue}{\textsf{[S]}}}
\newcommand{\Rmark}{\textcolor{sourceblue}{\textsf{[R]}}}
\newcommand{\Nmark}{\textcolor{newgreen}{\textsf{[N]}}}
\newcommand{\Cmark}{\textcolor{conjecturered}{\textsf{[C]}}}

\newenvironment{boxedremark}{%
  \par\smallskip\noindent\begin{center}\begin{minipage}{0.94\textwidth}%
  \color{black}\hrule\vspace{0.6em}\small
}{%
  \vspace{0.6em}\hrule\end{minipage}\end{center}\smallskip
}
\newenvironment{proofidea}{\par\smallskip\noindent\textit{Proof architecture.}\ }{\hfill$\square$\par\smallskip}
\lstdefinestyle{pseudo}{
  basicstyle=\ttfamily\small,
  backgroundcolor=\color{shadegray},
  frame=single,
  rulecolor=\color{rulegray},
  columns=fullflexible,
  keepspaces=true,
  showstringspaces=false,
  xleftmargin=0.7em,
  xrightmargin=0.7em,
  aboveskip=0.8em,
  belowskip=0.8em
}

\title{\bfseries Atomic Functions Beyond the Critical Dyadic Case\\[0.35em]
\large English reconstruction of Rvachev's chapter, repository synthesis, and new fractal--probabilistic frontiers}
\author{Research report prepared in the notation and house style of the \texttt{ProveIt} Fabius corpus}
\date{28 August 2026}

\begin{document}
\maketitle

\begin{abstract}
This report reconstructs in English, corrects at the level permitted by the surviving page images, and substantially expands the attached Russian material on Rvachev's atomic functions.  The source begins with compactly supported solutions of differential--functional equations, reduces their existence to cancellation of zeros in an infinite Fourier product, specializes to the one-parameter densities
\[
 \widehat h_a(t)=\prod_{j\ge1}\sinc(t a^{-j}),\qquad a>1,
\]
and treats the critical member $h_2=\Up$, polynomial reproduction by dyadic translates, and the first auxiliary $\operatorname{Fup}$ function.  The exposition below preserves that organization and notation while separating four proof statuses: translated source results, established repository/literature connections, new theorems proved here, and explicit conjectures.

The principal new contribution is a complete geometric description of the separated regime $a>2$.  The set on which $h_a$ fails to be locally polynomial is exactly the two-map Cantor attractor
\[
 K_a=\left\{\sum_{j\ge1}\eps_j a^{-j}:\eps_j\in\{-1,1\}\right\}.
\]
Every complementary gap of generation $n$ carries a polynomial of exact degree $n$.  This gives exact $L^1$ and $L^\infty$ norms of all derivatives, the Hausdorff dimension $\log2/\log a$, a closed geometric zeta function and its complete lattice of complex dimensions, an exact tube formula with nonconstant logarithmic oscillation, and a geometric law for the local polynomial degree.  As $a\downarrow2$, the rescaled degree converges to an exponential random variable, making the dyadic nowhere-analytic function a precise critical limit of full-measure polynomial splines.

A second new layer gives Bernoulli-number cumulants and complete Bell-polynomial moments for every $a>1$, a Gaussian limit as $a\downarrow1$, a uniform-law limit as $a\to\infty$, and an exact general-base negative-Laplace decomposition.  The latter contains a genuinely one-periodic correction whose nonzero Fourier coefficients are
\[
 -\frac{\Gamma(-\chi_k)\zeta(1-\chi_k)}{\log a},\qquad
 \chi_k=\frac{2\pi\ii k}{\log a},
\]
thereby extending the dyadic Gamma--zeta correction to arbitrary base.  A lower-$W_{-1}$ saddle coordinate and an all-orders periodically modulated endpoint program are formulated.  Reproducible Python code, numerical audits, figures, and data tables accompany the report.
\end{abstract}

\begin{boxedremark}
\textbf{Proof-status legend.}
\Smark\ denotes a statement translated or cleanly reconstructed from the attached source;
\Rmark\ denotes a connection already established in the surrounding repository or classical literature;
\Nmark\ denotes a result proved in this report; and
\Cmark\ denotes a plausible but presently unproved conjecture or research program.  A statement may carry more than one marker when a source observation is sharpened here.
\end{boxedremark}

\tableofcontents
\newpage

\section{Scope, sources, and normalization}
\label{sec:scope}

\subsection{What was translated and how}

The attached \texttt{rvachev.tex} is an OCR transcription of a Russian chapter titled
\emph{Atomic Functions}.  It contains five sections: the general definition and Fourier
existence criterion, the first-order two-shift equation, the special function $\Up$, polynomial
representation by translates of $\Up$, and the beginning of a section on functions
$\operatorname{Fup}_n$.  The attached PDF contains the surviving printed pages corresponding
mainly to the latter part of the first-order analysis and the $\Up$ section.  The OCR is often
reliable in prose and badly damaged in superscripts, factorials, binomial coefficients, and
binary subscripts.  Accordingly, this report follows three rules.

\begin{enumerate}
\item When the printed page makes a formula unambiguous, the formula is silently typeset in
      standard notation and the source equation number is retained in the discussion.
\item When the OCR and page image are structurally clear but individual indices are damaged,
      the invariant statement is reconstructed from the neighboring equations and verified
      independently.
\item When the surviving material does not determine every symbol, no fictitious exact
      transcription is supplied.  Instead, a clean equivalent formulation is given and the
      uncertainty is recorded in \cref{app:ocr}.
\end{enumerate}

The repository was audited recursively using its current document topology and its own
provenance maps.  The live corpus consists of the mathematical glossary, the primary exposition
\path{Fabius_Function_and_Rvachev_Up/Fabius_Function_and_Rvachev_Up.tex}, the
non-elementarity paper, the Lean walkthrough, the canonical frontier volume, the thematically
organized frontier drafts, and the vendored paper sources.  The frozen \path{archive/} was
cross-checked through its document-by-document manifest, which records the original location,
first committed version, and whether a living copy supersedes it.  The draft \path{MANIFEST.md}
was used as the global index for the reorganized 28 August 2026 frontier tree, including the
Fourier-decay corpus and the consolidated Thue--Morse, $q$-series, transform, inverse, sampling,
representation, spectral-arithmetic, and multi-topic volumes.  Current canonical TeX sources
were read directly; superseded historical repetitions were used for provenance and
nonduplication checks rather than counted as independent mathematical evidence.

The primary document already contains the formalized dyadic probability, Fourier, arithmetic,
$q$-binomial, inverse, and lower-Lambert theories.  The present report therefore concentrates on
the translated \emph{general-base} family $h_a$ and on deductions that were not found as exact
statements in that corpus.  In particular, the report does not repackage the existing dyadic
formula atlases, Fourier-decay audits, transform compilations, or lower-Lambert coefficient
machinery.  This is a nonduplication audit, not a claim of priority over the entire literature.

\subsection{Fourier convention}

Throughout,
\begin{equation}
 \widehat f(t)=\int_{\R}\e^{\ii tx}f(x)\dd x,
 \qquad
 f(x)=\frac1{2\pi}\int_{\R}\e^{-\ii tx}\widehat f(t)\dd t,
 \label{eq:fourier-convention}
\end{equation}
and
\[
 \sinc z=\begin{cases}\sin z/z,&z\ne0,\\1,&z=0.\end{cases}
\]
All logarithms are natural.  Binary digits are indexed from the least significant digit:
$p_j(m)\in\{0,1\}$ denotes the coefficient of $2^{j-1}$ in $m$, and
$\wt(m)=\sum_jp_j(m)$.

\subsection{Three related dyadic functions}

The repository carefully separates the bounded Fabius distribution function $F$, the even
compactly supported bump $\Up$, and the signed global extension.  Only the first two are
needed here.  Their bridge is
\begin{equation}
 \Up(x)=F(1-|x|),\qquad F(x)=\Up(x-1)\quad(0\le x\le1).
 \label{eq:F-up-bridge}
\end{equation}
Thus the left-endpoint behavior of $\Up$ is exactly the small-argument behavior of $F$.  The
general family $h_a$ below contains $\Up$ at $a=2$ but does not, for arbitrary $a$, inherit all
special dyadic identities of $F$.

\section{Atomic functions and the Fourier-product criterion}
\label{sec:general-atomic}

\subsection{The motivating compromise: local and universal}

\Smark\ The source begins from a practical tension in approximation theory.  Algebraic and
trigonometric polynomials are universal in the sense that the same class approximates
functions of arbitrarily high smoothness, but they are nonlocal.  Classical splines are local,
and their compact support yields sparse or banded numerical matrices, but a fixed spline
degree is not universal: increasing smoothness requires increasing degree.  Rvachev's
proposal is to use compactly supported $C^\infty$ functions obtained as infinite convolutions
of interval indicators.  Their Fourier transforms are products
\begin{equation}
 \widehat\varphi(t)=\prod_{k=1}^{\infty}\sinc(\alpha_k t),
 \qquad \alpha_k>0,
 \qquad \sum_{k\ge1}\alpha_k<\infty.
 \label{eq:general-sinc-product}
\end{equation}
Finite convolutions give ordinary splines.  Under the source's regularity hypotheses---in
particular for the geometric scales used below---the infinite convolution gives a compactly
supported infinitely differentiable limit.  Mere summability of an arbitrary sequence
$(\alpha_k)$ controls the support but, by itself, should not be read as a complete $C^\infty$
criterion.  The geometrically decreasing choice $\alpha_k=2^{-k}$ produces $\Up$.

\subsection{Differential--functional equations}

\begin{definition}[Atomic function, source normalization]
\label{def:atomic}
An \emph{atomic function} is a nonzero compactly supported solution of
\begin{equation}
 P(D)y(x)=\lambda\sum_{k=1}^{M}c_k y(ax-b_k),
 \qquad a>1,
 \qquad D=-\ii\frac{\dd}{\dd x},
 \label{eq:atomic-equation}
\end{equation}
where $P$ is a polynomial and the right-hand side is a finite dilation--translation operator.
\end{definition}

Put
\begin{equation}
 \mathcal E(t)=\frac1a\sum_{k=1}^{M}c_k\e^{\ii tb_k/a}.
 \label{eq:E-mask}
\end{equation}
Taking Fourier transforms in \eqref{eq:atomic-equation} gives the exact scalar recursion
\begin{equation}
 P(t)Y(t)=\lambda\mathcal E(t)Y(t/a).
 \tag{S3.2}
 \label{eq:fourier-recursion-general}
\end{equation}
Iteration yields
\begin{equation}
 Y(t)\prod_{j=0}^{n}P(ta^{-j})
 =\lambda^{n+1}\prod_{j=0}^{n}\mathcal E(ta^{-j})Y(ta^{-n-1}).
 \tag{S3.7}
 \label{eq:iterated-general}
\end{equation}
This is the basic reason Fourier analysis is as effective for atomic equations as it is for
constant-coefficient ODEs.

\subsection{Zero cancellation and normalized solutions}

\Smark\ Suppose first that $Y(0)=\int y=1$.  Then the zero of lowest order of $P$ at the
origin must match the zero of the same order of $\mathcal E$, and the normalization fixes
$\lambda$.  Away from the origin, every zero introduced by the denominator in the infinite
iteration must be canceled by a scaled zero of the exponential polynomial $\mathcal E$.
The source phrases this as an injection between the multisets of zeros of $P$ and
$\mathcal E$ along multiplicative $a$-orbits.

\begin{theorem}[Fourier-product existence criterion, cleaned source statement]
\label{thm:source-existence}
Within the source class and its stated growth hypotheses, let $A$ be the multiset of zeros
of $P$ and $B$ the multiset of zeros of $\mathcal E$.  Assume the origin multiplicities and
normalization are compatible.  A normalized nonzero compactly supported solution of
\eqref{eq:atomic-equation} exists if and only if the zeros of
$P$ can be assigned, with multiplicity, to distinct zeros of $\mathcal E$ of the form
$a^{-m}\alpha$, $m\in\N_0$, along their dilation orbits.  Under this cancellation condition,
\begin{equation}
 Y(t)=\prod_{j=0}^{\infty}
       \frac{\lambda\mathcal E(ta^{-j})}{P(ta^{-j})},
 \tag{S3.8}
 \label{eq:source-product-general}
\end{equation}
with removable singularities filled in, and $y$ is recovered by Fourier inversion.
\end{theorem}

\begin{proofidea}
Necessity follows from \eqref{eq:iterated-general}: for fixed $t$, the terminal factor tends
to $Y(0)=1$, so an uncanceled denominator zero is impossible.  For sufficiency, the
cancellation makes \eqref{eq:source-product-general} entire.  The factors approach $1$
uniformly on compact sets, giving local uniform convergence.  The exponential type is
controlled by the summable translation scales, while repeated use of the functional equation
gives decay faster than every power on the real axis.  Paley--Wiener then gives compact
support, and Fourier inversion verifies the equation.
\end{proofidea}

\begin{warning}
The OCR preserves the zero-orbit criterion and product, but not every side condition from the
preceding pages with equal reliability.  Accordingly, \cref{thm:source-existence} records the
source mechanism inside its original analytic class; it is not an assertion that an arbitrary
polynomial and exponential mask satisfying only the displayed zero assignment automatically
produce a $C^\infty$ compactly supported inverse transform.
\end{warning}

\Rmark\ In modern language this theorem belongs to the same broad family as refinement and
two-scale equations: existence is encoded by a mask, dilation, and cancellation of spectral
obstructions.  The Daubechies--Lagarias theory studies finite masks and local regularity by
matrix products; Rvachev's equation differs in that a differential operator appears on the
left, but the Fourier recursion has the same renormalization structure
\cite{daubechieslagarias1,daubechieslagarias2}.

\subsection{The derivative ladder when the mass vanishes}

\Smark\ If $Y(0)=0$ and the zero has order $r$, write $Y(t)=t^rG(t)$ with $G(0)=1$.
The normalized function corresponding to $G$ solves the same cancellation problem with the
eigenvalue rescaled by $a^{-r}$, while the original compactly supported solution is a constant
multiple of the $r$th derivative of that normalized solution.  Consequently the admissible
eigenvalues form a geometric ladder and the derivative solutions are linearly independent.
This observation foreshadows the exact derivative formulas for $h_a$ and $\Up$.

\section{The first-order family $h_a$}
\label{sec:ha-definition}

\subsection{Reduction of the two-shift equation}

\Smark\ For $P(t)=-\ii t+\alpha$ and two shifts, translation and scaling reduce the equation
to
\begin{equation}
 y'(x)+\alpha y(x)=\lambda\bigl(y(ax-1)+c\,y(ax+1)\bigr).
 \tag{S3.22}
 \label{eq:first-order-normal}
\end{equation}
When $\alpha\ne0$, the zero-cancellation condition selects a countable family of coefficients
$c$; the source gives them explicitly through logarithms of the two-term exponential mask.
The central case for atomic-function theory is $\alpha=0$, where normalization forces
$c=-1$ and $\lambda=-a^2/2$.

\begin{definition}[Generalized Rvachev density]
\label{def:ha}
For $a>1$, define
\begin{equation}
 \Phi_a(t)=\prod_{j=1}^{\infty}\sinc(ta^{-j}),
 \qquad
 h_a(x)=\frac1{2\pi}\int_{\R}\e^{-\ii tx}\Phi_a(t)\dd t.
 \tag{S3.24--S3.25}
 \label{eq:ha-Fourier}
\end{equation}
\end{definition}

The Fourier recursion is
\begin{equation}
 \Phi_a(t)=\sinc(t/a)\Phi_a(t/a),
 \tag{S3.28}
 \label{eq:ha-fourier-refinement}
\end{equation}
and differentiation of the corresponding integral refinement gives
\begin{equation}
 h_a'(x)=\frac{a^2}{2}\bigl(h_a(ax+1)-h_a(ax-1)\bigr)
       =-\frac{a^2}{2}\bigl(h_a(ax-1)-h_a(ax+1)\bigr).
 \tag{S3.23}
 \label{eq:ha-diffeq}
\end{equation}
Equivalently,
\begin{equation}
 h_a(x)=\frac a2\int_{ax-1}^{ax+1}h_a(u)\dd u.
 \label{eq:ha-integral-refinement}
\end{equation}
The last formula immediately gives $0\le h_a(x)\le a/2$ once the probabilistic
interpretation is known.

\subsection{Infinite convolution and support}

\begin{theorem}[Probability representation, source]
\label{thm:probability-ha}
Let $U_1,U_2,\ldots$ be independent random variables uniformly distributed on $[-1,1]$ and
put
\begin{equation}
 X_a=\sum_{j=1}^{\infty}a^{-j}U_j.
 \tag{S3.26}
 \label{eq:Xa}
\end{equation}
Then $\Phi_a$ is the characteristic function of $X_a$, $h_a$ is its probability density,
and
\begin{equation}
 \supp h_a=[-b,b],\qquad b=\sum_{j\ge1}a^{-j}=\frac1{a-1}.
 \tag{S3.27}
 \label{eq:ha-support}
\end{equation}
The density is even, nonnegative, $C^\infty$, and has total mass $1$.
\end{theorem}

\begin{proof}
The characteristic function of $a^{-j}U_j$ is $\sinc(ta^{-j})$, so independence gives the
product.  Almost surely $|X_a|\le b$, and every point of $[-b,b]$ is approximated by finite
choices of the coordinates, so the topological support is the entire interval.  Along the
real axis the product decreases faster than every inverse power.  Indeed, if
$N\asymp\log_a|t|$, then the first $N$ factors are bounded by $a^j/|t|$, giving
\begin{equation}
 |\Phi_a(t)|\le
 \exp\!\left(-\frac{(\log|t|)^2}{2\log a}+\bigO(\log|t|)\right).
 \label{eq:lognormal-fourier-bound}
\end{equation}
Thus every $t^m\Phi_a(t)$ is integrable and $h_a\in C^\infty$.  Compact support also follows
from the exponential type $\sum a^{-j}=b$ by Paley--Wiener.
\end{proof}

\Rmark\ This is an infinite convolution in the sense of Jessen and Wintner
\cite{jessen1935}.  It is simultaneously a fixed point of a smoothing transform: $X_a$ has
the same law as $a^{-1}(U+X_a')$, where $X_a'$ is an independent copy.  The resulting
integral refinement \eqref{eq:ha-integral-refinement} is the continuous-mask analogue of a
subdivision equation.

\subsection{Finite sinc products as nonuniform box splines}

For $N\ge1$, set
\begin{equation}
 \Phi_{a,N}(t)=\prod_{j=1}^{N}\sinc(ta^{-j}),
 \qquad
 X_{a,N}=\sum_{j=1}^{N}a^{-j}U_j,
 \label{eq:finite-products}
\end{equation}
and let $h_{a,N}$ be the density of $X_{a,N}$.  It is a compactly supported piecewise
polynomial of degree $N-1$, a one-dimensional box spline with nonuniform directions
$a^{-1},\ldots,a^{-N}$.  If $w_j=a^{-j}$, $W_N=\sum_{j\le N}w_j$, and
$W_S=\sum_{j\in S}w_j$, then inclusion--exclusion gives
\begin{equation}
 h_{a,N}(x)=\frac1{2^N(N-1)!\prod_{j=1}^{N}w_j}
 \sum_{S\subseteq\{1,\ldots,N\}}(-1)^{|S|}
 \bigl(x+W_N-2W_S\bigr)_+^{N-1}.
 \label{eq:finite-box-spline}
\end{equation}
The knots are the signed sums $-W_N+2W_S$.

\begin{proposition}[Finite-product convergence]
\label{prop:finite-convergence}
For the coupling using the same coordinates,
\begin{equation}
 |X_a-X_{a,N}|\le\frac{a^{-N}}{a-1}
 \quad\text{almost surely},
 \qquad
 \Wass(\mathcal L(X_a),\mathcal L(X_{a,N}))
 \le\frac{a^{-N}}{a-1}.
 \label{eq:Wasserstein-finite}
\end{equation}
Moreover, for every fixed $m\ge0$, $h_{a,N}\to h_a$ in $C^m(\R)$ as $N\to\infty$.
On bounded $t$-sets,
\begin{equation}
 |\Phi_a(t)-\Phi_{a,N}(t)|
 \le \frac{t^2a^{-2(N+1)}}{6(1-a^{-2})}.
 \label{eq:finite-product-rate}
\end{equation}
\end{proposition}

\begin{proof}
The tail bound is the geometric series.  The product estimate follows by telescoping and
$|1-\sinc z|\le z^2/6$.  For $C^m$ convergence, choose $N_0>m+2$.  For $N\ge N_0$ both
$\Phi_{a,N}$ and $\Phi_a$ are dominated by the integrable function
$\prod_{j=1}^{N_0}|\sinc(ta^{-j})|$ after multiplication by $|t|^m$.  Dominated convergence
and Fourier inversion finish the proof.
\end{proof}

\section{Taylor coefficients, moments, Bernoulli numbers, and Bell polynomials}
\label{sec:moments}

\subsection{The source recurrence}

Write
\begin{equation}
 \Phi_a(t)=\sum_{k=0}^{\infty}c_{2k}(a)t^{2k},
 \qquad c_0(a)=1.
 \label{eq:coefficient-expansion}
\end{equation}
Odd coefficients vanish by symmetry.  Comparing coefficients in
\eqref{eq:ha-fourier-refinement} gives the source recurrence
\begin{equation}
 c_{2k}(a)=\frac1{a^{2k}-1}
 \sum_{\ell=0}^{k-1}
 \frac{(-1)^{k-\ell}c_{2\ell}(a)}{(2k-2\ell+1)!},
 \qquad k\ge1.
 \tag{S3.32}
 \label{eq:source-c-recurrence}
\end{equation}
The even moments
\begin{equation}
 \mu_{2k}(a)=\int_{-b}^{b}x^{2k}h_a(x)\dd x
            =(-1)^k(2k)!c_{2k}(a)
 \tag{S3.36--S3.37}
 \label{eq:moment-coefficient}
\end{equation}
therefore satisfy an exact rational recurrence whenever $a$ is rational or algebraic.
The source iterates \eqref{eq:source-c-recurrence} to obtain superfactorial decay estimates
for $c_{2k}(a)$; these estimates underlie the very rapid endpoint algorithms in the dyadic
case.

\subsection{A closed cumulant formula}

\Nmark\ The recurrence becomes more transparent after passing from moments to cumulants.

\begin{theorem}[Bernoulli cumulants]
\label{thm:cumulants}
All odd cumulants of $X_a$ vanish, and for $m\ge1$,
\begin{equation}
 \boxed{\quad
 \kappa_{2m}(X_a)=
 \frac{2^{2m}B_{2m}}{2m\,(a^{2m}-1)}.
 \quad}
 \label{eq:cumulant-formula}
\end{equation}
Consequently
\begin{equation}
 \mu_n(a)=\mathcal B_n\bigl(0,\kappa_2,0,\kappa_4,\ldots,\kappa_n\bigr),
 \label{eq:Bell-moments}
\end{equation}
where $\mathcal B_n$ is the complete exponential Bell polynomial.
\end{theorem}

\begin{proof}
For $U\sim\operatorname{Unif}[-1,1]$,
\[
 \log\E\e^{zU}=\log\frac{\sinh z}{z}
 =\sum_{m\ge1}\frac{2^{2m}B_{2m}}{2m(2m)!}z^{2m}.
\]
Cumulants add under independence and scale homogeneously, hence
\[
 \kappa_{2m}(X_a)=\kappa_{2m}(U)\sum_{j\ge1}a^{-2mj}
 =\frac{2^{2m}B_{2m}}{2m}\frac1{a^{2m}-1}.
\]
The Bell-polynomial identity is the standard moment--cumulant relation.
\end{proof}

The first values are
\begin{align}
 \kappa_2&=\frac1{3(a^2-1)},
 &\kappa_4&=-\frac2{15(a^4-1)},
 &\kappa_6&=\frac{16}{63(a^6-1)},\label{eq:first-cumulants}\\
 \mu_2&=\kappa_2,
 &\mu_4&=3\kappa_2^2+\kappa_4,
 &\mu_6&=15\kappa_2^3+15\kappa_2\kappa_4+\kappa_6.
 \label{eq:first-moments}
\end{align}
This is computationally preferable to repeatedly expanding the infinite product.

\subsection{$q$-Lambert structure}

Let $q=a^{-2}$.  Then
\begin{equation}
 \frac1{a^{2m}-1}=\frac{q^m}{1-q^m},
 \label{eq:Lambert-factor}
\end{equation}
a classical Lambert-series factor.  The logarithm of the characteristic function becomes
\begin{equation}
 \log\Phi_a(t)=
 \sum_{m\ge1}
 \frac{(-1)^m2^{2m}B_{2m}}{2m(2m)!}
 \frac{q^m}{1-q^m}t^{2m}.
 \label{eq:logPhi-Lambert}
\end{equation}
Expanding $q^m/(1-q^m)=\sum_{r\ge1}q^{mr}$ gives a divisor-polynomial expansion
\begin{equation}
 \log\Phi_a(t)=\sum_{n\ge1}q^n
 \sum_{m\mid n}
 \frac{(-1)^m2^{2m}B_{2m}}{2m(2m)!}t^{2m}.
 \label{eq:divisor-polynomial}
\end{equation}
This is the direct bridge from the general $a$ family to the repository's $q$-series and
$q$-binomial material.  At $a=2$, $q=1/4$, while the half-base arithmetic of the bounded
Fabius function introduces the complementary base $1/2$ through dyadic scaling.

\section{The separated regime $a>2$: smooth fractal splines}
\label{sec:separated}

\subsection{The central plateau and source polynomial pieces}

Put
\begin{equation}
 b=\frac1{a-1},\qquad
 b_1=b\left(1-\frac2a\right)=\frac{a-2}{a(a-1)}.
 \label{eq:b-b1}
\end{equation}
For $a>2$, the two intervals on which $h_a(ax\pm1)$ can be nonzero are disjoint.
Equation \eqref{eq:ha-diffeq} therefore gives $h_a'(x)=0$ on $[-b_1,b_1]$.
Using \eqref{eq:ha-integral-refinement} and total mass $1$ gives the source plateau
\begin{equation}
 h_a(x)=\frac a2\qquad(|x|\le b_1).
 \tag{S3.40}
 \label{eq:central-plateau}
\end{equation}
The source then observes recursively that $h_a$ is linear on two intervals, quadratic on
four intervals, and so forth.  The union of all such polynomial intervals has full measure;
the remaining null set is Cantor-like.  We now identify this geometry exactly.

\subsection{The singular attractor and its gaps}

Define the contractions
\begin{equation}
 S_-(x)=\frac{x-1}{a},\qquad S_+(x)=\frac{x+1}{a}.
 \label{eq:IFS-maps}
\end{equation}

\begin{definition}[Address Cantor set]
\label{def:Ka}
For $a>2$, let
\begin{equation}
 K_a=S_-(K_a)\cup S_+(K_a)
 =\left\{\sum_{j=1}^{\infty}\eps_j a^{-j}:\eps_j\in\{-1,1\}\right\}.
 \label{eq:Ka}
\end{equation}
The central gap is $G_0=(-b_1,b_1)$.  For a word
$w=(\eps_1,\ldots,\eps_n)\in\{-,+\}^n$, write
$S_w=S_{\eps_1}\circ\cdots\circ S_{\eps_n}$ and
$G_w=S_w(G_0)$.
\end{definition}

\begin{theorem}[Exact polynomial locus]
\label{thm:polynomial-locus}
For $a>2$,
\begin{equation}
 [-b,b]\setminus K_a=\bigsqcup_{n\ge0}\ \bigsqcup_{|w|=n}G_w.
 \label{eq:gap-decomposition}
\end{equation}
On every generation-$n$ gap $G_w$, the function $h_a$ is a polynomial of exact degree $n$.
Consequently the set of points admitting a neighborhood on which $h_a$ is polynomial is
precisely $[-b,b]\setminus K_a$.
\end{theorem}

\begin{proof}
The first-level images $S_\pm([-b,b])$ are disjoint and leave exactly the central gap $G_0$.
Standard iteration of a separated two-map IFS gives \eqref{eq:gap-decomposition}.

Differentiate \eqref{eq:ha-diffeq} repeatedly.  With
\begin{equation}
 \delta_k=(-1)^{\wt(k-1)},\qquad 1\le k\le2^n,
 \label{eq:delta-TM}
\end{equation}
one obtains
\begin{equation}
 h_a^{(n)}(x)=
 \left(\frac{a^2}{2}\right)^n a^{n(n-1)/2}
 \sum_{k=1}^{2^n}\delta_k\,
 h_a\!\left(a^nx+\sum_{j=1}^{n}a^{j-1}(-1)^{p_j(k-1)}\right).
 \tag{S3.49}
 \label{eq:ha-derivative}
\end{equation}
For $x$ in a fixed generation-$n$ gap, exactly one summand in
\eqref{eq:ha-derivative} lands in the central plateau and every other summand lands outside
$[-b,b]$.  Therefore $h_a^{(n)}$ is a nonzero constant on that gap.  Differentiating once
more gives zero there because the surviving plateau is constant.  Hence the degree is
exactly $n$.

Every point of $K_a$ is accumulated by gaps of arbitrarily large generation.  A single
polynomial neighborhood of degree $d$ cannot contain a generation-$n$ gap with $n>d$.
Thus no point of $K_a$ belongs to the polynomial locus.
\end{proof}

\begin{corollary}[Exact nonanalytic locus]
\label{cor:nonanalytic-locus}
Within its support, $h_a$ is real analytic exactly on $[-b,b]\setminus K_a$, where it is
locally a polynomial.  It is nonanalytic at every point of $K_a$; on
$\R\setminus[-b,b]$ it is identically zero and therefore analytic.
\end{corollary}

\begin{proof}
The complement is polynomial by \cref{thm:polynomial-locus}.  If $h_a$ were analytic in a
neighborhood of a point of $K_a$, that neighborhood would contain a gap.  On the gap the
analytic function agrees with a polynomial, so the identity theorem would make it that same
polynomial throughout the neighborhood, contradicting \cref{thm:polynomial-locus}.
\end{proof}

This sharpens the source's statement that $h_a$ is polynomial on a full-measure set and
nonanalytic on a Cantor-type null set.

\subsection{Breakpoints and exact moment evaluation}

\Smark\ The source describes the boundary points created after $n$ differentiations as
\begin{equation}
 x_0=a^{-n}\left(\pm b+\sum_{k=0}^{n-1}\pm a^k\right).
 \tag{S3.50}
 \label{eq:source-breakpoints}
\end{equation}
These are exactly the endpoints of the cylinder intervals $S_w([-b,b])$ and of their
complementary gaps.  At such points, the Taylor remainder formula based at $-b$,
\begin{equation}
 h_a(x)=\frac1{n!}\int_{-b}^{x}h_a^{(n+1)}(t)(x-t)^n\dd t,
 \label{eq:Taylor-endpoint}
\end{equation}
combined with \eqref{eq:ha-derivative}, reduces values and derivatives to finite linear
combinations of the even moments \eqref{eq:moment-coefficient}.  The OCR of the source's
fully indexed formulas (S3.54)--(S3.58) is too damaged to reproduce character by character,
but their invariant content is exactly this finite moment reduction.  A stable computational
version is given in \cref{alg:gap-polynomial}.

\begin{algorithm}[Exact polynomial on a prescribed gap]
\label{alg:gap-polynomial}
Given $a>2$ and a word $w$ of length $n$:
\begin{enumerate}
\item compute the left endpoint $x_w$ of $G_w$ from the affine maps $S_\pm$;
\item compute moments $\mu_{0},\mu_2,\ldots,\mu_{2\floor{(n-1)/2}}$ by
      \cref{thm:cumulants};
\item use \eqref{eq:Taylor-endpoint} and \eqref{eq:ha-derivative} to evaluate
      $h_a^{(r)}(x_w)$ for $0\le r\le n$;
\item return
\[
 P_w(x)=\sum_{r=0}^{n}\frac{h_a^{(r)}(x_w)}{r!}(x-x_w)^r.
\]
By \cref{thm:polynomial-locus}, $P_w=h_a$ throughout $G_w$.
\end{enumerate}
\end{algorithm}

\subsection{Exact derivative norms}

The source gives the exact dyadic norm law for $\Up$.  The disjoint-support geometry shows
that it extends to every $a\ge2$.

\begin{theorem}[Exact derivative norms]
\label{thm:derivative-norms}
For every $a\ge2$ and $n\ge0$,
\begin{equation}
 \boxed{\quad
 \|h_a^{(n)}\|_{\infty}
 =\frac{a^{(n+1)(n+2)/2}}{2^{n+1}},
 \qquad
 \|h_a^{(n)}\|_{1}=a^{n(n+1)/2}.
 \quad}
 \label{eq:general-derivative-norms}
\end{equation}
\end{theorem}

\begin{proof}
The integral refinement gives $\|h_a\|_\infty\le a/2$, with equality on the plateau for
$a>2$ and at $x=0$ for $a=2$.  The $2^n$ support intervals of the summands in
\eqref{eq:ha-derivative} are disjoint when $a>2$ and meet only at endpoints when $a=2$.
Thus there is no cancellation in either the supremum or the absolute integral.  The common
coefficient is $(a^2/2)^n a^{n(n-1)/2}$.  Multiplying it by $a/2$ gives the first formula.
For the second, each rescaled copy has integral $a^{-n}$, so
\[
 \|h_a^{(n)}\|_1=
 \left(\frac{a^2}{2}\right)^n a^{n(n-1)/2}\,2^n a^{-n}
 =a^{n(n+1)/2}.
\]
\end{proof}

At $a=2$, \eqref{eq:general-derivative-norms} reduces to the source identity
\begin{equation}
 \|\Up^{(n)}\|_\infty=\|\Up^{(n)}\|_1=2^{n(n+1)/2}.
 \tag{S3.64}
\end{equation}

\section{Fractal geometry and the critical transition}
\label{sec:fractal}

\subsection{Dimension and geometric zeta function}

The central gap has length
\begin{equation}
 \ell_0=2b_1=\frac{2(a-2)}{a(a-1)},
 \label{eq:ell0}
\end{equation}
and the generation-$n$ gaps have length $\ell_n=\ell_0a^{-n}$ with multiplicity $2^n$.

\begin{theorem}[Dimension and complex dimensions]
\label{thm:complex-dimensions}
For $a>2$,
\begin{equation}
 \dim_HK_a=\dim_MK_a=D_a:=\frac{\log2}{\log a}.
 \label{eq:dimension-Ka}
\end{equation}
The geometric zeta function of the complementary fractal string is
\begin{equation}
 \zeta_{\mathcal L_a}(s)
 =\sum_{n\ge0}2^n\ell_n^s
 =\frac{\ell_0^s}{1-2a^{-s}},
 \qquad \Re s>D_a.
 \label{eq:fractal-zeta}
\end{equation}
It extends meromorphically to $\C$ with simple poles
\begin{equation}
 s_k=D_a+\frac{2\pi\ii k}{\log a},\qquad k\in\Z,
 \label{eq:complex-dimensions}
\end{equation}
and
\begin{equation}
 \operatorname*{res}_{s=s_k}\zeta_{\mathcal L_a}(s)
 =\frac{\ell_0^{s_k}}{\log a}.
 \label{eq:zeta-residue}
\end{equation}
\end{theorem}

\begin{proof}
The open-set condition holds because $a>2$, and the similarity equation
$2a^{-D}=1$ gives \eqref{eq:dimension-Ka} by Hutchinson's theorem
\cite{hutchinson1981}.  The gap sum is geometric.  Solving $1-2a^{-s}=0$ gives the pole
lattice, and differentiation of the denominator gives residue $\ell_0^{s_k}/\log a$.
\end{proof}

\Rmark\ The vertical arithmetic progression \eqref{eq:complex-dimensions} is the canonical
lattice pattern of a self-similar fractal string in the sense of Lapidus and van
Frankenhuijsen \cite{lapidus2003,lapidusbook}.  Here it enters the analysis of a $C^\infty$
function: the singular set of the smooth spline carries the same complex dimensions as its
complementary gap string.

\subsection{An exact tube formula}

For $\eps>0$, define the inner tube volume
\begin{equation}
 V_a(\eps)=
 \left|\left\{x\in[-b,b]\setminus K_a:\dist(x,K_a)<\eps\right\}\right|.
 \label{eq:tube-def}
\end{equation}
Each gap of length $\ell$ contributes $\min(2\eps,\ell)$.

\begin{theorem}[Exact tube formula]
\label{thm:tube-formula}
Let $a>2$ and $0<2\eps<\ell_0$.  Put
\begin{equation}
 N(\eps)=\floor{\log_a\frac{\ell_0}{2\eps}}.
 \label{eq:N-eps}
\end{equation}
Then
\begin{equation}
 \boxed{\quad
 V_a(\eps)=2\eps\bigl(2^{N(\eps)+1}-1\bigr)
 +\frac{\ell_0(2/a)^{N(\eps)+1}}{1-2/a}.
 \quad}
 \label{eq:exact-tube}
\end{equation}
If $x=\log_a(\ell_0/(2\eps))$ and $\theta=\fract{x}$, then
\begin{equation}
 \eps^{D_a-1}V_a(\eps)=\mathcal M_a(\theta)-2\eps^{D_a},
 \label{eq:tube-normalized}
\end{equation}
where the one-periodic profile is
\begin{equation}
 \mathcal M_a(\theta)=\ell_0^{D_a}
 \left[
 2^{2-D_a-\theta}
 +\frac{2^{1-D_a}a^{(D_a-1)(1-\theta)}}{1-2/a}
 \right],\qquad0\le\theta<1.
 \label{eq:Minkowski-profile}
\end{equation}
The profile is continuous under periodic identification and nonconstant.  Hence $K_a$ is not
Minkowski measurable, although its upper and lower Minkowski contents are finite and positive.
\end{theorem}

\begin{proof}
The gaps with $n\le N$ are wider than $2\eps$ and contribute $2\eps$ each; the remaining
gaps are completely covered.  Summing the finite geometric count and the infinite geometric
length tail gives \eqref{eq:exact-tube}.  Substitution of
$\eps=(\ell_0/2)a^{-x}$ and $a^{D_a}=2$ gives
\eqref{eq:tube-normalized}--\eqref{eq:Minkowski-profile}.  The two positive exponentials in
$\theta$ cannot sum to a constant for $a>2$.
\end{proof}

The logarithmic average of the oscillatory content is explicit:
\begin{equation}
 \int_0^1\mathcal M_a(\theta)\dd\theta
 =\ell_0^{D_a}2^{1-D_a}
 \left(\frac1{\log2}+\frac1{\log(a/2)}\right).
 \label{eq:average-Minkowski}
\end{equation}

\subsection{The local polynomial degree as a random variable}

Choose a point uniformly with respect to normalized Lebesgue measure on $[-b,b]$.
Since $K_a$ has measure zero, the point lies in a unique gap.  Let $N_a$ be the exact degree
of the polynomial piece containing it.

\begin{theorem}[Geometric degree law]
\label{thm:degree-law}
For $a>2$,
\begin{equation}
 \Prob(N_a=n)=\frac{a-2}{a}\left(\frac2a\right)^n,
 \qquad n=0,1,2,\ldots.
 \label{eq:degree-geometric}
\end{equation}
Therefore
\begin{equation}
 \E N_a=\frac2{a-2},
 \qquad
 \Var(N_a)=\frac{2a}{(a-2)^2}.
 \label{eq:degree-moments}
\end{equation}
Let $p_a=(a-2)/a$.  As $a\downarrow2$,
\begin{equation}
 p_aN_a\ \xrightarrow{\ d\ }\ \operatorname{Exp}(1),
 \qquad
 \frac{a-2}{2}N_a\ \xrightarrow{\ d\ }\ \operatorname{Exp}(1).
 \label{eq:degree-critical-limit}
\end{equation}
\end{theorem}

\begin{proof}
The total length of all generation-$n$ gaps is $2^n\ell_0a^{-n}$.  Dividing by
$2b=2/(a-1)$ gives \eqref{eq:degree-geometric}.  The moment formulas are those of a
geometric law on $\N_0$.  Finally,
\[
 \Prob(p_aN_a>x)=(1-p_a)^{\floor{x/p_a}+1}\longrightarrow\e^{-x}.
\]
\end{proof}

This theorem makes the phrase ``polynomial almost everywhere'' quantitative.  Near the
critical base, a typical polynomial piece has a very large degree, of order $(a-2)^{-1}$.

\subsection{The dyadic phase transition}

As $a\downarrow2$,
\begin{equation}
 D_a=1-\frac{a-2}{2\log2}+\bigO((a-2)^2),
 \qquad
 \ell_0=(a-2)+\bigO((a-2)^2).
 \label{eq:critical-expansions}
\end{equation}
Thus the null singular set thickens to dimension $1$, the central plateau collapses, and the
mean local degree diverges.  At $a=2$ the two IFS pieces merely touch; the complement has no
open gaps at all.  The limit function is nowhere analytic on its support, as the source states.
This is a genuine regularity phase transition, not merely a change in support geometry.

\section{The critical member $h_2=\Up$}
\label{sec:up}

\subsection{Fourier product, equation, support, and derivatives}

\Smark\ At $a=2$,
\begin{equation}
 \Up(x)=\frac1{2\pi}\int_{\R}\e^{-\ii tx}
 \prod_{k=1}^{\infty}\sinc(t2^{-k})\dd t,
 \tag{S3.60}
 \label{eq:up-product}
\end{equation}
\begin{equation}
 \Up'(x)=2\Up(2x+1)-2\Up(2x-1),
 \qquad \supp\Up=[-1,1].
 \tag{S3.61--S3.62}
 \label{eq:up-equation}
\end{equation}
Since
\[
 \sum_{j=1}^{n}2^{j-1}(-1)^{p_j(k-1)}=2^n+1-2k,
\]
formula \eqref{eq:ha-derivative} becomes
\begin{equation}
 \Up^{(n)}(x)=2^{n(n+1)/2}
 \sum_{k=1}^{2^n}\delta_k\Up(2^nx+2^n+1-2k).
 \tag{S3.63}
 \label{eq:up-derivative}
\end{equation}
The exact norm identity is the $a=2$ case of \cref{thm:derivative-norms}.

\subsection{Taylor polynomials and nowhere analyticity}

\Smark\ At a dyadic rational represented at depth $n$, the Taylor series of $\Up$
terminates after degree at most $n$; at every nondyadic point of $[-1,1]$, the Taylor series has
radius of convergence zero.
Hence $\Up$ is nonanalytic at every point of its support.  The dyadic points are not analytic
exceptions: their finite Taylor polynomials do not agree with $\Up$ on any neighborhood.
This is the critical counterpart of \cref{cor:nonanalytic-locus}: for $a>2$ the polynomial
Taylor series is locally exact off $K_a$, while for $a=2$ there are no open polynomial gaps.

\subsection{Moment formulas and exact dyadic values}

Let
\begin{equation}
 \mu_{2r}=\int_{-1}^{1}x^{2r}\Up(x)\dd x,
 \qquad
 m_p^+=\int_0^1x^p\Up(x)\dd x.
 \label{eq:up-moments}
\end{equation}
Then
\begin{equation}
 \mu_{2r}=(-1)^r(2r)!c_{2r}(2),
 \tag{S3.65}
 \label{eq:up-even-moment}
\end{equation}
and integration by parts in \eqref{eq:up-equation} gives, for $p\ge0$,
\begin{equation}
 m_p^+=\frac1{(p+1)2^{p+1}}
 \sum_{r=0}^{\floor{(p+1)/2}}
 \binom{p+1}{2r}\mu_{2r}.
 \label{eq:half-moment-recurrence}
\end{equation}
The especially useful endpoint value is
\begin{equation}
 \Up(-1+2^{-n})=
 \frac{m_{n-1}^+}{2^{n(n-1)/2}(n-1)!},
 \qquad n\ge1.
 \tag{S3.68}
 \label{eq:up-endpoint-dyadic}
\end{equation}
More generally, for $1\le k\le2^n$,
\begin{align}
 \Up(-1+k2^{-n})
 &=\frac{2^{-n(n-1)/2}}{n!}
 \sum_{j=1}^{k}\delta_j
 \int_{-1}^{1}\Up(t)
 \left(k-j+\frac12-\frac t2\right)^n\dd t \notag\\
 &=\frac{2^{-n(n-1)/2}}{n!}
 \sum_{j=1}^{k}\delta_j
 \sum_{r=0}^{\floor{n/2}}
 \binom n{2r}\left(k-j+\frac12\right)^{n-2r}2^{-2r}\mu_{2r}.
 \label{eq:up-general-dyadic}
\end{align}
This is a cleaned version of the source's formula (S3.67).

\subsection{Binary recursion and a rapidly convergent evaluation scheme}

For $n\ge1$, define
\begin{equation}
 \varphi_n(x)=\Up(x)+\Up(x-2^{-n}).
 \label{eq:phi-n}
\end{equation}
On
$I_n=[-1+2^{-n},-1+2^{-n+1}]$, the source observes that
$\varphi_n^{(n+1)}=0$, because the first two Thue--Morse signs cancel.  Therefore
\begin{equation}
 \varphi_n(x)=\sum_{r=0}^{n}
 \frac{\Up^{(r)}(-1+2^{-n})}{r!}(x+1-2^{-n})^r,
 \qquad x\in I_n,
 \tag{S3.70}
 \label{eq:phi-polynomial}
\end{equation}
and
\begin{equation}
 \Up(x)=\sum_{r=0}^{n}
 \frac{\Up^{(r)}(-1+2^{-n})}{r!}(x+1-2^{-n})^r
 -\Up(x-2^{-n}).
 \tag{S3.71}
 \label{eq:up-recursion}
\end{equation}
Iterating \eqref{eq:up-recursion} according to the binary digits of $x+1$ produces the
source's series (S3.73).  The exact printed indices are partly corrupted in the OCR, but the
algorithm and its error bound are unambiguous:
\begin{equation}
 |R_N(x)|\le\Up(-1+2^{-N}),
 \label{eq:binary-remainder}
\end{equation}
and the right side decreases faster than any geometric sequence because of
\eqref{eq:up-endpoint-dyadic} and the superfactorial moment bounds.  This is the conceptual
ancestor of modern exact dyadic evaluators in the repository and in \cite{haugland2016}.

\section{Thue--Morse signs and polynomial reproduction}
\label{sec:thue-morse}

\subsection{The derivative operator is a Prouhet filter}

The signs in the source recurrence
\begin{equation}
 \delta_{2k-1}=\delta_k,\qquad \delta_{2k}=-\delta_k,
 \qquad\delta_1=1,
 \label{eq:delta-recursion}
\end{equation}
are exactly
$\delta_k=(-1)^{\wt(k-1)}$.  Their generating polynomial is
\begin{equation}
 \sum_{k=0}^{2^n-1}(-1)^{\wt(k)}z^k
 =\prod_{j=0}^{n-1}(1-z^{2^j}).
 \label{eq:TM-product}
\end{equation}
Thus the $n$th derivative formula is the product of $n$ two-branch difference filters.  The
zero of order $n$ at $z=1$ gives the Prouhet cancellation
\begin{equation}
 \sum_{k=0}^{2^n-1}(-1)^{\wt(k)}k^r=0,
 \qquad 0\le r<n.
 \label{eq:Prouhet}
\end{equation}
This is why high-order polynomial pieces and exact dyadic identities emerge from what first
looks like a functional-differential equation.

\subsection{Characteristic polynomial for translate coefficients}

\Smark\ The source considers locally finite sums
\begin{equation}
 \phi(x)=\sum_{k\in\Z}c_k\Up(x-k2^{-n}).
 \tag{S3.76}
 \label{eq:translate-sum}
\end{equation}
For $\phi$ to be a polynomial of degree at most $n$, its $(n+1)$st derivative must vanish,
which is equivalent to the constant-coefficient difference system
\begin{equation}
 \sum_{j=1}^{2^{n+1}}c_{k+j}\delta_j=0,
 \qquad k\in\Z.
 \tag{S3.77}
 \label{eq:coefficient-recurrence}
\end{equation}
The characteristic polynomial is
\begin{equation}
 Q_n(z)=\sum_{k=1}^{2^{n+1}}\delta_kz^{2^{n+1}-k}.
 \tag{S3.79}
 \label{eq:Qn-def}
\end{equation}
The source recurrence
\begin{equation}
 Q_n(z)=Q_{n-1}(z)(z^{2^n}-1),\qquad Q_0(z)=z-1,
 \tag{S3.80--S3.81}
 \label{eq:Qn-recurrence}
\end{equation}
gives
\begin{equation}
 \boxed{\quad Q_n(z)=\prod_{j=0}^{n}(z^{2^j}-1).\quad}
 \label{eq:Qn-product}
\end{equation}

A root of exact order $2^m$ has multiplicity $n-m+1$ for $1\le m\le n$, while $1$ has
multiplicity $n+1$.  Hence every solution of \eqref{eq:coefficient-recurrence} is a finite sum
of polynomial--exponential sequences
\begin{equation}
 c_k=\sum_{\omega:Q_n(\omega)=0}
 \left(\sum_{r=0}^{\operatorname{mult}(\omega)-1}A_{\omega,r}k^r\right)\omega^k.
 \label{eq:recurrence-general-solution}
\end{equation}
The high multiplicity of $1$ is the polynomial-reproduction mechanism; the other roots of
unity describe oscillatory null modes.

\subsection{Uniqueness of the universal representing function}

\Smark\ The source proves a uniqueness theorem strengthening earlier Rvachev results.
A clean statement is as follows.

\begin{theorem}[Source uniqueness principle]
\label{thm:source-uniqueness}
Let $\varphi\in C_c^\infty(\R)$ have support $[-1,1]$ and mass $1$.  Assume
\begin{equation}
 \lim_{n\to\infty}
 \frac{\|\varphi^{(n)}\|_\infty}{2^{(n+1)(n+2)/2}}=0,
 \label{eq:uniqueness-growth}
\end{equation}
and assume that for every $n$ there are coefficients $c_k^{(n)}$ such that the locally finite
sum
\begin{equation}
 \sum_{k\in\Z}c_k^{(n)}\varphi(x-k2^{-n})=x^n.
 \label{eq:universal-reproduction}
\end{equation}
Then $\varphi=\Up$.
\end{theorem}

The proof builds a sequence $\varphi_n$ satisfying the same mass and reproduction conditions
and
\begin{equation}
 \varphi_n'(x)=2\varphi_{n+1}(2x+1)-2\varphi_{n+1}(2x-1).
 \label{eq:source-phi-chain}
\end{equation}
Comparing with the identical integral recursion for $\Up$, the derivative-growth hypothesis
forces the successive differences to contract to zero.  The growth condition is sharp in the
source's sense: $\Up+\alpha\Up'$ retains mass and polynomial reproduction but violates
\eqref{eq:uniqueness-growth} with limiting ratio $|\alpha|$.

\section{The $\operatorname{Fup}$ construction and a natural hierarchy}
\label{sec:fup}

\subsection{The surviving source definition}

\Smark\ The attached source ends shortly after using
\begin{equation}
 \sinc(t/2)=\cos(t/4)\sinc(t/4)
 \label{eq:sinc-halving}
\end{equation}
to rewrite
\begin{equation}
 \widehat\Up(t)=\cos(t/4)\left(\sinc(t/4)\right)^2
 \prod_{j=3}^{\infty}\sinc(t2^{-j}).
 \tag{S3.98}
 \label{eq:up-Fup-factor}
\end{equation}
It defines $\operatorname{Fup}_1$ by
\begin{equation}
 \widehat{\operatorname{Fup}_1}(t)=
 \left(\sinc(t/4)\right)^2
 \prod_{j=3}^{\infty}\sinc(t2^{-j}),
 \tag{S3.99}
 \label{eq:Fup1}
\end{equation}
and obtains
\begin{equation}
 \Up(x)=\frac12\left(
 \operatorname{Fup}_1(x-1/4)+\operatorname{Fup}_1(x+1/4)
 \right).
 \tag{S3.100}
 \label{eq:Fup1-average}
\end{equation}
No later pages of the $\operatorname{Fup}_n$ section are present in the attachment.

\subsection{A canonical continuation}

\Nmark\ There is a simple hierarchy forced by the same factorization.  Put
\begin{equation}
 H(t)=\prod_{j=2}^{\infty}\sinc(t2^{-j})
 \label{eq:H-tail}
\end{equation}
and define $G_n$ for $n\ge0$ by
\begin{equation}
 \widehat G_n(t)=H(t)\sinc(t2^{-(n+1)}).
 \label{eq:Gn-hat}
\end{equation}
Then $G_0=\Up$ and $G_1=\operatorname{Fup}_1$.

\begin{proposition}[Dyadic Fup ladder]
\label{prop:Fup-ladder}
For $n\ge1$,
\begin{equation}
 G_{n-1}(x)=\frac12\left(
 G_n(x-2^{-(n+1)})+G_n(x+2^{-(n+1)})
 \right).
 \label{eq:Gn-recursion}
\end{equation}
Moreover,
\begin{equation}
 G_n\longrightarrow 2\Up(2\,\cdot)
 \quad\text{in }C^m(\R)\text{ for every fixed }m.
 \label{eq:Gn-limit}
\end{equation}
The support radius of $G_n$ is $1/2+2^{-(n+1)}$.
\end{proposition}

\begin{proof}
Identity \eqref{eq:sinc-halving} gives
$\widehat G_{n-1}(t)=\cos(t2^{-(n+1)})\widehat G_n(t)$, which is
\eqref{eq:Gn-recursion}.  Pointwise,
$\widehat G_n(t)\to H(t)$, and $H$ is the Fourier transform of the density of
$\sum_{j\ge2}2^{-j}U_j=X_2/2$, namely $2\Up(2x)$.  Since
$|\widehat G_n-H|\le2|H|$ and $|t|^mH(t)\in L^1$, dominated Fourier inversion gives
$C^m$ convergence.  The support radius is the sum of the uniform half-widths in
\eqref{eq:Gn-hat}.
\end{proof}

Iterating \eqref{eq:Gn-recursion} expresses $\Up$ as an average of $2^n$ translates of
$G_n$.  In the limit, the translate locations become the usual centered dyadic Bernoulli
sum, whose law is uniform on $[-1/2,1/2]$.  The limit recovers the first-step convolution
identity for $\Up$.  This makes the $\operatorname{Fup}$ construction a finite Rademacher
factorization of the uniform smoothing mask.

\section{Parameter limits: Gaussian, critical, and uniform}
\label{sec:parameter-limits}

\subsection{The many-small-summands limit $a\downarrow1$}

The variance is
\begin{equation}
 \sigma_a^2=\Var(X_a)=\frac1{3(a^2-1)}.
 \label{eq:variance-a}
\end{equation}
Let $Z_a=X_a/\sigma_a$.

\begin{theorem}[Gaussian limit and first characteristic correction]
\label{thm:gaussian-limit}
As $a\downarrow1$,
\begin{equation}
 Z_a\xrightarrow{\ d\ }N(0,1).
 \label{eq:Gaussian-limit}
\end{equation}
For every fixed $T>0$, uniformly for $|t|\le T$,
\begin{equation}
 \E\e^{\ii tZ_a}
 =\e^{-t^2/2}\left[
 1-\frac{a^2-1}{20(a^2+1)}t^4
 +\bigO_T((a-1)^2)
 \right].
 \label{eq:Gaussian-first-correction}
\end{equation}
More generally, the standardized $2m$th cumulant is
\begin{equation}
 \lambda_{2m}(a)=
 \frac{3^m2^{2m}B_{2m}(a^2-1)^m}{2m(a^{2m}-1)}
 =\bigO((a-1)^{m-1}),\qquad m\ge2.
 \label{eq:standardized-cumulant}
\end{equation}
\end{theorem}

\begin{proof}
The triangular-array coefficients in
$Z_a=\sum_j\sqrt{3(a^2-1)}a^{-j}U_j$ have maximal magnitude tending to zero and total
variance one, so the Lindeberg theorem applies.  Alternatively, insert
\eqref{eq:standardized-cumulant} in the logarithm of the characteristic function.  The exact
fourth standardized cumulant is
\begin{equation}
 \lambda_4(a)=-\frac65\frac{a^2-1}{a^2+1}.
 \label{eq:lambda4}
\end{equation}
All higher contributions to the logarithm are $\bigO_T((a-1)^2)$, and exponentiation gives
\eqref{eq:Gaussian-first-correction}.
\end{proof}

\Cmark\ A natural next theorem is a uniform local Edgeworth expansion for the normalized
densities and their derivatives.  The characteristic expansion and superpolynomial Fourier
decay make this plausible, but uniform control as $a\downarrow1$ requires a careful split
between central and large frequencies.

\subsection{The single-dominant-summand limit $a\to\infty$}

Put $Y_a=aX_a$.  Then
\begin{equation}
 Y_a=U_1+T_a,
 \qquad
 |T_a|\le\frac1{a-1}.
 \label{eq:large-a-coupling}
\end{equation}

\begin{theorem}[Uniform limit with quantitative bounds]
\label{thm:uniform-limit}
Let $U\sim\operatorname{Unif}[-1,1]$.  For $a>2$,
\begin{equation}
 \Wass(\mathcal L(Y_a),\mathcal L(U))\le\frac1{a-1},
 \qquad
 \TV(\mathcal L(Y_a),\mathcal L(U))\le\frac1{2(a-1)}.
 \label{eq:uniform-bounds}
\end{equation}
If $g_a$ is the density of $Y_a$, then
\begin{equation}
 g_a(y)=\frac12
 \quad\text{for}\quad
 |y|\le\frac{a-2}{a-1}.
 \label{eq:uniform-core}
\end{equation}
Thus the convergence is exact on an expanding central interval; all discrepancy is confined
to two boundary layers, each of width $2/(a-1)$ (total width $4/(a-1)$).
\end{theorem}

\begin{proof}
The coupling \eqref{eq:large-a-coupling} gives the Wasserstein bound.  Conditional on
$T_a=r$, the law is uniform on $[r-1,r+1]$, whose total variation distance from
$[-1,1]$ is $|r|/2$; convexity of total variation gives the second bound.  Finally,
$g_a(y)=h_a(y/a)/a$, and the plateau \eqref{eq:central-plateau} becomes
\eqref{eq:uniform-core}.
\end{proof}

\section{Exact negative-Laplace renormalization in arbitrary base}
\label{sec:laplace-general}

\subsection{Endpoint shift and product}

Let
\begin{equation}
 Y_a=X_a+b\in[0,2b],
 \qquad f_a(y)=h_a(-b+y).
 \label{eq:endpoint-shift}
\end{equation}
Since $U_j+1$ is uniform on $[0,2]$,
\begin{equation}
 \int_0^{2b}\e^{-sy}f_a(y)\dd y
 =\prod_{j=1}^{\infty}
 \frac{1-\e^{-2sa^{-j}}}{2sa^{-j}}.
 \label{eq:endpoint-Laplace}
\end{equation}
Define, for $u>0$,
\begin{equation}
 \kappa(u)=\log\frac{1-\e^{-u}}{u},
 \qquad
 \Lambda_a(u)=\sum_{j=1}^{\infty}\kappa(ua^{-j}).
 \label{eq:Lambda-a}
\end{equation}
Then the logarithm of \eqref{eq:endpoint-Laplace} is $\Lambda_a(2s)$, and
\begin{equation}
 \Lambda_a(au)=\Lambda_a(u)+\kappa(u).
 \label{eq:Lambda-dilation}
\end{equation}

\subsection{The exact periodic decomposition}

Put $L_a=\log a$ and, for $x\in\R$,
\begin{equation}
 R_a(x)=\sum_{n=0}^{\infty}\log\bigl(1-\e^{-a^{x+n}}\bigr),
 \label{eq:Rax}
\end{equation}
which is negative and exponentially convergent.  Define
\begin{equation}
 P_a(x)=\Lambda_a(a^x)+\frac{L_a}{2}x^2-\frac{L_a}{2}x+R_a(x).
 \label{eq:Pa-def}
\end{equation}

\begin{theorem}[General-base exact periodic decomposition]
\label{thm:general-periodic}
For every $a>1$, $P_a$ is one-periodic.  For $u>0$,
\begin{equation}
 \boxed{\quad
 \Lambda_a(u)=
 -\frac{(\log u)^2}{2L_a}+\frac{\log u}{2}
 +P_a(\log_a u)+E_a(u),
 \quad}
 \label{eq:general-exact-Laplace}
\end{equation}
where
\begin{equation}
 E_a(u)=-R_a(\log_a u)
 =\sum_{n=0}^{\infty}-\log(1-\e^{-ua^n})\ge0.
 \label{eq:Ea}
\end{equation}
Moreover,
\begin{equation}
 0\le E_a(u)\le
 \frac{\e^{-u}}{(1-\e^{-u})(1-\e^{-(a-1)u})}.
 \label{eq:Ea-bound}
\end{equation}
In particular $E_a(u)=\bigO_a(\e^{-u})$ as $u\to\infty$.
\end{theorem}

\begin{proof}
Using \eqref{eq:Lambda-dilation} and
$R_a(x+1)=R_a(x)-\log(1-\e^{-a^x})$, one obtains
\[
 P_a(x+1)-P_a(x)
 =\kappa(a^x)+L_ax+R_a(x+1)-R_a(x)=0.
\]
Rearranging \eqref{eq:Pa-def} gives \eqref{eq:general-exact-Laplace}.  For the tail, use
$-\log(1-z)\le z/(1-z)$, $a^n\ge1+n(a-1)$, and sum the resulting geometric series.
\end{proof}

For $a=2$, \eqref{eq:Ea-bound} becomes
$E_2(u)\le\e^{-u}/(1-\e^{-u})^2$, exactly the dyadic form appearing in the primary
repository exposition.  Thus \cref{thm:general-periodic} is a base-$a$ renormalization of the
same mechanism rather than an unrelated analogy.

\subsection{Gamma--zeta Fourier coefficients}

Let
\begin{equation}
 \overline P_a=\int_0^1P_a(x)\dd x,
 \qquad \Psi_a(x)=P_a(x)-\overline P_a,
 \qquad
 \chi_k=\frac{2\pi\ii k}{L_a}.
 \label{eq:Psi-a}
\end{equation}

\begin{theorem}[General-base Gamma--zeta modes]
\label{thm:general-Fourier-modes}
For $k\ne0$,
\begin{equation}
 \boxed{\quad
 \widehat\Psi_a(k)=\int_0^1\Psi_a(x)\e^{-2\pi\ii kx}\dd x
 =-\frac{\Gamma(-\chi_k)\zeta(1-\chi_k)}{L_a}.
 \quad}
 \label{eq:general-Gamma-zeta}
\end{equation}
The Fourier series is absolutely convergent with all derivatives; in fact the coefficients
decay exponentially in $|k|$.  Therefore $P_a$ is real analytic and nonconstant.
\end{theorem}

\begin{proofidea}
In the strip $-1<\Re s<0$, integration by parts gives the Mellin transform
\begin{equation}
 \int_0^\infty\kappa(u)u^{s-1}\dd u
 =-\Gamma(s)\zeta(s+1).
 \label{eq:Mellin-kappa}
\end{equation}
Mellin summation of the dilation equation samples this transform at
$s=-\chi_k$, producing \eqref{eq:general-Gamma-zeta}; the polynomial terms in
\eqref{eq:general-exact-Laplace} affect only the zero mode.  Stirling's formula for
$\Gamma(\sigma+\ii t)$ gives exponential decay.  Nonconstancy follows already from the
nonzero first mode; more generally $\Gamma$ has no zeros and the classical zero-free line
$\zeta(1+\ii t)\ne0$ shows that every nonzero mode is nonzero.
\end{proofidea}

\Rmark\ Formula \eqref{eq:general-Gamma-zeta} is the arbitrary-base counterpart of the exact
dyadic periodic correction in the repository.  It also belongs to the general Mellin theory
of harmonic sums developed by Flajolet, Gourdon, and Dumas \cite{flajolet1995}.  The same
``lattice in logarithmic scale'' mechanism appears independently in
\cref{thm:complex-dimensions}; one periodicity comes from endpoint Laplace renormalization,
the other from the geometric gap string.

\subsection{Lower Lambert coordinate and the endpoint program}

Ignoring the bounded periodic derivative and exponentially small tail in the saddle equation
for \eqref{eq:endpoint-Laplace}, put
\begin{equation}
 r_a(y)a^{-r_a(y)}=\frac{y\sqrt a}{2}.
 \label{eq:r-Lambert-equation}
\end{equation}
The large solution is
\begin{equation}
 r_a(y)=-\frac1{\log a}
 W_{-1}\!\left(-\frac{y\sqrt a\log a}{2}\right),
 \qquad y\downarrow0.
 \label{eq:r-Lambert}
\end{equation}
This is the natural general-base lower-Lambert coordinate.  The exact saddle is shifted by a
bounded periodic term involving $P_a'$.

\begin{conjecture}[All-orders general-base endpoint expansion]
\label{conj:endpoint-all-orders}
Let $f_a(y)=h_a(-b+y)$.  For fixed $a>1$, the exact positive saddle $s_a(y)$ defined by
\begin{equation}
 y+2\Lambda_a'(2s_a(y))=0
 \label{eq:exact-saddle}
\end{equation}
exists and is unique for sufficiently small $y>0$.  There are real-analytic one-periodic
functions $A_{a,j}$ such that the saddle expansion
\begin{equation}
 f_a(y)\sim
 \frac{\exp\!\bigl(s_a(y)y+\Lambda_a(2s_a(y))\bigr)}
      {\sqrt{8\pi\Lambda_a''(2s_a(y))}}
 \left(1+\sum_{j\ge1}\frac{A_{a,j}(\log_a(2s_a(y)))}{r_a(y)^j}\right)
 \label{eq:general-saddle-conjecture}
\end{equation}
holds to arbitrary order.  Equivalently, $\log f_a(y)$ has a transseries whose leading scale
is $-(\log a)r_a(y)^2/2$, whose next terms include a linear term and
$-\tfrac12\log r_a(y)$, and whose order-one oscillation is
$P_a(r_a(y)+\tfrac12)$ after normalization.
\end{conjecture}

At $a=2$, normalization and a half-unit phase shift reduce this program to the exact
lower-$W_{-1}$ and Gamma--zeta hierarchy already developed in the repository.  The point of
\cref{conj:endpoint-all-orders} is not to rederive that dyadic case, but to identify the
correct general-base structure and the periodic term that any sharp formula must retain.

\section{Connections to existing theories}
\label{sec:connections}

\subsection{Infinite convolutions and smoothing transforms}

The probability law of $X_a$ is an infinite convolution of compactly supported absolutely
continuous measures.  Unlike Bernoulli convolutions, whose regularity may depend delicately on
arithmetic properties of the contraction, every digit here is uniform and each convolution
step adds smoothing.  Nevertheless, the \emph{address geometry} of the separated regime is
the same two-map Cantor geometry.  This separation between a smooth full-support density and
a fractal nonanalytic locus is one of the distinctive features of $h_a$.

\subsection{Refinement equations, subdivision, and wavelets}

The Fourier equation
$\Phi_a(t)=\sinc(t/a)\Phi_a(t/a)$ is a refinement relation with an analytic, nonfinite mask.
The finite approximants $h_{a,N}$ are box splines, while the limit is infinitely smooth.
Daubechies--Lagarias refinement theory emphasizes finite masks, matrix products, and Hölder
regularity \cite{daubechieslagarias1,daubechieslagarias2}; the present family supplies a
complementary exactly solvable model in which local regularity is not merely estimated but
encoded by a geometric gap address.

\subsection{Thue--Morse, Prouhet cancellation, and roots of unity}

The derivative signs, dyadic value formulas, and polynomial-reproduction recurrence are all
manifestations of the same factorization \eqref{eq:TM-product}.  The source's characteristic
polynomial \eqref{eq:Qn-product} records not only the zero at $1$ responsible for Prouhet
cancellation but the full hierarchy of $2$-power roots of unity and their multiplicities.
This explains the simultaneous appearance of polynomial coefficients, trigonometric
sequences, and $2$-adic valuations in Rvachev's representation theorem.

\subsection{Bernoulli and Bell polynomials}

Bernoulli numbers enter through the logarithm of $\sinh z/z$; Bell polynomials reconstruct
moments from cumulants.  In the dyadic corpus, Bernoulli and Bell polynomials also arise in
Euler--Maclaurin and endpoint expansions.  Here their role is more elementary and more
general: they give a closed moment calculus for every $a>1$ before any dyadic specialization.

\subsection{Two distinct logarithmic lattices}

Two periodic phenomena occur in this report.
\begin{enumerate}
\item The gap string has complex dimensions
      $D_a+2\pi\ii k/\log a$, producing tube oscillations in
      $\log_a(1/\eps)$.
\item The negative-Laplace equation has Fourier modes indexed by the same vertical spacing,
      producing endpoint oscillations in $\log_a u$.
\end{enumerate}
They have different amplitudes and arise from different objects, but both are consequences
of an exact multiplicative lattice.  A promising future problem is to determine whether
there is a direct transform carrying one periodic profile into the other.

\section{Numerical experiments and reproducibility}
\label{sec:numerics}

The accompanying script \path{experiments.py} uses only NumPy, Matplotlib, and mpmath.  It
creates all figures and CSV tables in the archive.  Random experiments use the fixed seed
\texttt{20260828}; the omitted random-series tail is deterministically below $10^{-12}$.

\subsection{Gap hierarchy and exact polynomial degree}

\begin{figure}[H]
\centering
\includegraphics[width=0.96\textwidth]{figures/gap_hierarchy_a_2_6.pdf}
\caption{The gaps $G_w$ for $a=2.6$.  Vertical level $n$ is both the gap generation and,
by \cref{thm:polynomial-locus}, the exact degree of $h_a$ on that interval.}
\label{fig:gap-hierarchy}
\end{figure}

\subsection{Tube oscillation}

\begin{figure}[H]
\centering
\includegraphics[width=0.96\textwidth]{figures/tube_oscillation_a_3.pdf}
\caption{The exact normalized tube volume for $a=3$.  The approach to a nonconstant
one-periodic profile verifies the lattice oscillation in \cref{thm:tube-formula}.}
\label{fig:tube-oscillation}
\end{figure}

\subsection{Critical local-degree limit}

\begin{figure}[H]
\centering
\includegraphics[width=0.96\textwidth]{figures/degree_critical_limit.pdf}
\caption{Distribution functions of $p_aN_a$ approaching $1-\e^{-x}$ as $a\downarrow2$.}
\label{fig:degree-limit}
\end{figure}

\subsection{General-base periodic correction}

\begin{figure}[H]
\centering
\includegraphics[width=0.96\textwidth]{figures/periodic_correction_bases.pdf}
\caption{Centered periodic corrections $P_a-\overline P_a$.  The dyadic oscillation is
small but nonzero; its amplitude grows substantially for the displayed larger bases.}
\label{fig:periodic-correction}
\end{figure}

On a grid, the maximum residual in $P_a(x+1)-P_a(x)$ is below
$3\times10^{-14}$ for $a\in\{1.2,1.5,2,2.6,3,5\}$.  Numerical Fourier coefficients agree
with \eqref{eq:general-Gamma-zeta} to between $10^{-12}$ and $10^{-16}$ in the supplied
table; the larger dyadic first-mode error reflects resolving a Fourier mode of amplitude about
$10^{-6}$ with double-precision quadrature, not a visible failure of the identity.

\subsection{Gaussian cumulant scaling}

\begin{figure}[H]
\centering
\includegraphics[width=0.96\textwidth]{figures/gaussian_cumulant_limit.pdf}
\caption{The exact standardized fourth cumulant and its first asymptotic term.}
\label{fig:gaussian-limit}
\end{figure}

For $a=2.6$, a fixed-seed simulation of $250{,}000$ random sums confirms the exact Bell
moments through order $10$ at the expected Monte Carlo accuracy.  Exact values and errors are
stored in \path{data/moment_validation_a_2_6.csv}.

\FloatBarrier

\section{Conjectures and future research directions}
\label{sec:future}

\begin{conjecture}[Overlap-regime nowhere analyticity]
\label{conj:overlap}
For every $1<a\le2$, $h_a$ is nonanalytic at every point of the interior of its support.
\end{conjecture}
The case $a=2$ is the source theorem.  For $1<a<2$ the two refinement branches overlap, so
the gap-degree argument disappears.  A proof may require a local Fourier or quasi-analyticity
criterion adapted to the dilation equation.

\begin{conjecture}[Critical profile of polynomial pieces]
\label{conj:critical-profile}
There is a nontrivial scaling limit for the polynomial containing a Lebesgue-typical point
when $a=2+\eta$, after rescaling its degree by $\eta/2$, its interval by its random gap length,
and its amplitude by the exact derivative norm in \cref{thm:derivative-norms}.
\end{conjecture}
The exponential degree law identifies the first marginal.  The remaining problem is a joint
limit of word addresses, polynomial coefficients, and moments as the gaps collapse.

\begin{conjecture}[All-order general-base saddle expansion]
The formal saddle program in \cref{conj:endpoint-all-orders} is valid uniformly for $a$ in
compact subsets of $(1,\infty)$, with coefficients generated by universal Bell polynomials in
the derivatives of $P_a$.
\end{conjecture}
This would connect the repository's dyadic lower-Lambert coefficients to a continuous base
parameter and make Bernoulli--Bell structure appear at both the moment and endpoint levels.

\begin{conjecture}[Arithmetic of algebraic breakpoints]
For algebraic $a>2$, exact values of $h_a$ and its derivatives at cylinder endpoints belong
to the field generated by $a$ and the even moments, with denominator growth controlled by
the ideals generated by $a^{2m}-1$.
\end{conjecture}
The source already proves rational dependence on $a$ at the formal level.  The cumulant
formula suggests a systematic algebraic-denominator theory generalizing the dyadic rational
arithmetic.

\begin{conjecture}[A bridge between the two periodic profiles]
There exists a natural spectral operator, built from the smoothing transform or its transfer
operator, relating the tube profile $\mathcal M_a$ of \eqref{eq:Minkowski-profile} to the
negative-Laplace profile $P_a$ of \eqref{eq:Pa-def}.
\end{conjecture}
The two profiles have the same logarithmic period but very different Fourier amplitudes.
Even a rigorous obstruction to a simple relation would clarify how geometry and endpoint
probability decouple.

\subsection*{Concrete research program}

The following projects appear especially promising.

\begin{enumerate}
\item Formalize the $a>2$ IFS, exact polynomial locus, derivative norms, degree law, and tube
      formula in Lean, using the repository's existing weighted-uniform distribution API.
\item Develop a transfer-operator proof or disproof of \cref{conj:overlap}.
\item Derive the first two general-base saddle coefficients explicitly and check their
      specialization against the formal dyadic coefficients.
\item Study total positivity, variation diminution, and approximation order of the finite
      box splines $h_{a,N}$ and their limit.
\item Extend the two-branch family to multiple uniform digits and nongeometric scales.  The
      separated singular locus then becomes a graph-directed or self-affine set, while the
      Fourier product remains explicit.
\item Complete the historical $\operatorname{Fup}_n$ translation from a full source copy and
      compare it with the canonical hierarchy of \cref{prop:Fup-ladder}.
\item Investigate stable algorithms for exact polynomial evaluation on deep gaps and for
      certified endpoint inversion via $W_{-1}$.
\item Quantify the convergence $h_{a,N}\to h_a$ in Sobolev, Besov, and weighted Fourier norms,
      and compare with subdivision regularity bounds.
\item Connect the roots-of-unity null modes in \eqref{eq:recurrence-general-solution} to
      wavelet vanishing moments and multiresolution decompositions generated by $\Up$.
\end{enumerate}

\section{Conclusion}

The attached chapter contains substantially more than the isolated definition of the
Rvachev bump.  It presents a general Fourier-product existence mechanism, a derivative
ladder, the one-parameter density $h_a$, a probability model, superfast moment recurrences,
full-measure polynomial splines for $a>2$, the critical nowhere-analytic function $\Up$,
exact dyadic values, polynomial reproduction, and a uniqueness principle.  In modern terms,
it sits at the intersection of infinite convolution, refinement equations, box splines,
automatic sequences, and compactly supported smooth analysis.

The general parameter $a$ reveals a particularly clean frontier.  For $a>2$, the smooth
function is locally polynomial off an exact Cantor set, and every geometric statistic of
that set---dimension, complex dimensions, tube volume, and degree distribution---has a closed
form.  At $a=2$ the gaps close, the typical degree diverges, and nowhere analyticity takes
over.  At the opposite parameter extremes, the law becomes Gaussian after normalization as
$a\downarrow1$ and uniform after scaling as $a\to\infty$.  Finally, the negative-Laplace
renormalization and Gamma--zeta periodic correction survive in every base, providing the
correct setting for a general lower-Lambert endpoint theory.

\appendix

\section{Repository corpus map and nonduplication protocol}
\label{app:repo-audit}

The repository tree consulted on 28 August 2026 separates living documents, temporary frontier
inboxes, vendored literature, and frozen historical versions.  The following map records how
those layers were used.

\begin{longtable}{@{}p{0.23\textwidth}p{0.31\textwidth}p{0.38\textwidth}@{}}
\toprule
Corpus layer & Representative contents & Role in this report\\
\midrule
\endhead
Living reference volumes & Glossary; primary Fabius--Rvachev exposition; non-elementarity
paper; Lean walkthrough & Fixed notation, theorem-status conventions, and the boundary between
formalized and exploratory claims.\\
Canonical frontier volume & \path{non-formalized-research-frontiers.tex} & Checked to avoid
repeating the established dyadic, $q$-series, inverse, transform, Thue--Morse, and endpoint
programs.\\
Frontier draft manifest & The Fourier-decay corpus and the consolidated exponents/$q$-series,
spectra/arithmetic, integration/transforms, inverse/sampling, representations, and collected
frontier volumes & Used as the recursive topic index after the 28 August 2026 reorganization;
exact-title overlap was treated as a warning to narrow or redirect a proposed result.\\
Vendored papers & Local TeX sources for the Fabius papers and related analytic literature &
Source-level comparison of normalizations, named results, and historical attributions.\\
Frozen archive & Early notes, first-generation primary drafts, supplements, $q$-calculus,
walkthroughs, glossaries, standalone studies, and former frontier dossiers & Provenance and
historical nonduplication only.  The archive manifest explicitly identifies living or
consolidated successors, so repeated text was not promoted as independent support.\\
Attached Russian chapter & OCR TeX plus the surviving printed PDF pages & Translation basis for
the atomic-function criterion, the family $h_a$, the spline geometry for $a>2$, the critical
function $\Up$, dyadic values, polynomial reproduction, and the opening of the $\operatorname{Fup}$
hierarchy.\\
\bottomrule
\end{longtable}

The nonduplication test was formula-level rather than title-level.  Before designating a result
\Nmark, its defining formula, parameter range, and conclusion were compared with the canonical
primary and frontier sources and with the draft manifest.  Results retained as new here are
therefore concentrated in the non-dyadic geometry of $h_a$: the exact Cantor nonanalytic locus,
exact gap degrees, derivative norms, fractal zeta and tube formulas, the critical degree law,
general-base cumulant limits, and the arbitrary-base Gamma--zeta periodic renormalization.
This procedure cannot establish historical priority beyond the audited repository and cited
literature; it does establish that the report does not intentionally duplicate an exact theorem
already indexed in the repository snapshot.

\section{OCR reconstruction ledger}
\label{app:ocr}

The following recurring OCR errors were normalized.

\begin{longtable}{@{}p{0.25\textwidth}p{0.29\textwidth}p{0.38\textwidth}@{}}
\toprule
OCR pattern & Reconstruction & Reason\\
\midrule
\endhead
\texttt{ta-j}, \texttt{a-i} & $ta^{-j}$, $a^{-j}$ & Geometric Fourier scales and the printed superscripts.\\
\texttt{l l}, \texttt{j j} & $l!$, $j!$ & Taylor denominators, confirmed by neighboring remainder formulas.\\
\texttt{C\_n\textasciicircum 2} in exponents & $\binom n2$ & Source notation for triangular exponents.\\
\texttt{up}, Cyrillic glyphs resembling ``ir, separated letters & $\Up$ & Cyrillic/Latin OCR mixing.\\
\texttt{Pj(k)} & $p_j(k)$ & The $j$th binary digit of $k$.\\
\texttt{delta} indices shifted by one & $\delta_k=(-1)^{\wt(k-1)}$ & Matches the printed recursion $\delta_{2k-1}=\delta_k$, $\delta_{2k}=-\delta_k$.\\
Damaged (S3.54)--(S3.58) & finite moment reduction at cylinder endpoints & The page image determines the structure but not every superscript reliably; no false literal transcription was attempted.\\
Damaged (S3.73)--(S3.75) & binary recursion plus error $\le\Up(-1+2^{-N})$ & The iterative identity and remainder are clear; the displayed digit-index series is not fully recoverable from OCR alone.\\
Uniqueness growth denominator & $2^{(n+1)(n+2)/2}$ & Confirmed by the source's counterexample $\Up+\alpha\Up'$ and its limiting ratio.\\
\bottomrule
\end{longtable}

\section{Notation concordance}
\label{app:notation}

\begin{longtable}{@{}p{0.18\textwidth}p{0.37\textwidth}p{0.37\textwidth}@{}}
\toprule
Symbol & Meaning here & Relation to source/repository\\
\midrule
\endhead
$h_a$ & density with Fourier product $\prod_{j\ge1}\sinc(ta^{-j})$ & Source notation; $h_2=\Up$.\\
$\Phi_a$ & characteristic function of $X_a$ & Source often writes $F_a(t)$; renamed to avoid conflict with bounded Fabius $F$.\\
$F$ & bounded Fabius function on $[0,1]$ with exterior constants & Primary repository convention.\\
$\Up$ & even compactly supported Rvachev bump & Source \texttt{up}.\\
$X_a$ & $\sum a^{-j}U_j$ & Source random variable $\xi(a)$.\\
$b$ & support radius $1/(a-1)$ & Source (S3.27).\\
$b_1$ & central plateau half-width $(a-2)/(a(a-1))$ & Source immediately before (S3.40).\\
$K_a$ & exact nonanalytic Cantor set for $a>2$ & New precise identification of the source's $\operatorname{sing}_a$.\\
$\delta_k$ & $(-1)^{\wt(k-1)}$ & Source recursive sign.\\
$c_{2m}(a)$ & Taylor coefficient of $\Phi_a$ & Source (S3.29)--(S3.32).\\
$\mu_{2m}(a)$ & even moment of $h_a$ & Source (S3.36)--(S3.38).\\
$P_a$ & general-base periodic Laplace correction & Generalizes the primary repository's dyadic correction.\\
\bottomrule
\end{longtable}

\section{Pseudocode for exact gap geometry}

\begin{lstlisting}[style=pseudo]
input: base a>2, maximum generation N
b1 = (a-2)/(a*(a-1))
gaps = [(-b1,b1)]
for n = 0,...,N:
    output every interval in gaps with degree n
    next = []
    for (left,right) in gaps:
        next.append(((left-1)/a,(right-1)/a))
        next.append(((left+1)/a,(right+1)/a))
    gaps = sort(next)
\end{lstlisting}

For moments, compute cumulants from \eqref{eq:cumulant-formula} and use the Bell recurrence
\begin{equation}
 \mu_0=1,
 \qquad
 \mu_n=\sum_{j=1}^{n}\binom{n-1}{j-1}\kappa_j\mu_{n-j}.
 \label{eq:Bell-recurrence}
\end{equation}

\section{Reproducibility manifest}
\label{app:reproducibility}

The archive contains:
\begin{itemize}
\item this LaTeX source and its compiled PDF;
\item \path{experiments.py}, with detailed comments and a fixed random seed;
\item vector and raster versions of all five figures;
\item CSV files for geometry and derivative norms, moment validation, periodicity residuals,
      and Gamma--zeta Fourier-mode validation;
\item \path{README.md}, \path{requirements.txt}, and \path{SHA256SUMS}.
\end{itemize}
Run
\begin{lstlisting}[style=pseudo]
python experiments.py --output .
latexmk -pdf Atomic_Functions_Beyond_Dyadic.tex
\end{lstlisting}
from the report directory.  The figures are deterministic; the Monte Carlo table is
reproducible because the seed is fixed.

\begin{thebibliography}{99}

\bibitem{jessen1935}
B.~Jessen and A.~Wintner,
\newblock Distribution functions and the Riemann zeta function,
\newblock \emph{Transactions of the American Mathematical Society} \textbf{38} (1935), 48--88.

\bibitem{fabius1966}
J.~Fabius,
\newblock A probabilistic example of a nowhere analytic $C^\infty$-function,
\newblock \emph{Zeitschrift f\"ur Wahrscheinlichkeitstheorie und Verwandte Gebiete}
\textbf{5} (1966), 173--174.

\bibitem{rvachev1971}
V.~L.~Rvachev and V.~A.~Rvachev,
\newblock A certain finite function,
\newblock \emph{Dopovidi Akademii Nauk Ukrains'koi RSR, Series A} (1971), 705--707, 764.

\bibitem{rvachev1979}
V.~L.~Rvachev and V.~A.~Rvachev,
\newblock \emph{Non-classical Methods of Approximation Theory in Boundary Value Problems},
\newblock Naukova Dumka, Kiev, 1979.  Russian.

\bibitem{rvachev1990}
V.~A.~Rvachev,
\newblock Compactly supported solutions of functional-differential equations and their applications,
\newblock \emph{Russian Mathematical Surveys} \textbf{45} (1990), no.~1, 87--120.

\bibitem{arias1982}
J.~Arias de Reyna,
\newblock Definici\'on y estudio de una funci\'on indefinidamente diferenciable de soporte compacto,
\newblock \emph{Revista de la Real Academia de Ciencias Exactas, F\'isicas y Naturales de Madrid}
\textbf{76} (1982), no.~1, 21--38; English translation,
\emph{An infinitely differentiable function with compact support: Definition and properties},
arXiv:1702.05442 (2017).

\bibitem{arias2018}
J.~Arias de Reyna,
\newblock Arithmetic of the Fabius function,
\newblock \emph{INTEGERS} \textbf{18} (2018), Article A51; arXiv:1702.06487.

\bibitem{haugland2016}
J.~K.~Haugland,
\newblock Evaluating the Fabius function,
\newblock arXiv:1609.07999 (2016).

\bibitem{allouche2003}
J.-P.~Allouche and J.~Shallit,
\newblock \emph{Automatic Sequences: Theory, Applications, Generalizations},
\newblock Cambridge University Press, 2003.

\bibitem{daubechieslagarias1}
I.~Daubechies and J.~C.~Lagarias,
\newblock Two-scale difference equations I: Existence and global regularity of solutions,
\newblock \emph{SIAM Journal on Mathematical Analysis} \textbf{22} (1991), 1388--1410.

\bibitem{daubechieslagarias2}
I.~Daubechies and J.~C.~Lagarias,
\newblock Two-scale difference equations II: Local regularity, infinite products of matrices and fractals,
\newblock \emph{SIAM Journal on Mathematical Analysis} \textbf{23} (1992), 1031--1079.

\bibitem{hutchinson1981}
J.~E.~Hutchinson,
\newblock Fractals and self-similarity,
\newblock \emph{Indiana University Mathematics Journal} \textbf{30} (1981), no.~5, 713--747.

\bibitem{lapidus2003}
M.~L.~Lapidus and M.~van Frankenhuijsen,
\newblock Complex dimensions of self-similar fractal strings and Diophantine approximation,
\newblock \emph{Experimental Mathematics} \textbf{12} (2003), no.~1, 41--69.

\bibitem{lapidusbook}
M.~L.~Lapidus and M.~van Frankenhuijsen,
\newblock \emph{Fractal Geometry, Complex Dimensions and Zeta Functions}, second edition,
\newblock Springer Monographs in Mathematics, Springer, 2013.

\bibitem{corless1996}
R.~M.~Corless, G.~H.~Gonnet, D.~E.~G.~Hare, D.~J.~Jeffrey, and D.~E.~Knuth,
\newblock On the Lambert $W$ function,
\newblock \emph{Advances in Computational Mathematics} \textbf{5} (1996), 329--359.

\bibitem{flajolet1995}
P.~Flajolet, X.~Gourdon, and P.~Dumas,
\newblock Mellin transforms and asymptotics: Harmonic sums,
\newblock \emph{Theoretical Computer Science} \textbf{144} (1995), 3--58.

\bibitem{andrews1976}
G.~E.~Andrews,
\newblock \emph{The Theory of Partitions},
\newblock Addison--Wesley, 1976.

\bibitem{gasper2004}
G.~Gasper and M.~Rahman,
\newblock \emph{Basic Hypergeometric Series}, second edition,
\newblock Cambridge University Press, 2004.

\bibitem{proveit}
V.~Reshetnikov,
\newblock \texttt{ProveIt}: formal and expository development of the Fabius--Rvachev system,
\newblock GitHub repository, \url{https://github.com/VladimirReshetnikov/ProveIt}, snapshot consulted 28 August 2026.

\end{thebibliography}

\end{document}
